\documentclass[11pt,a4paper]{article}
\usepackage[margin=1in]{geometry}
\usepackage{amsmath,amssymb,bm}
\usepackage{hyperref}
\usepackage{graphicx}
\usepackage{booktabs}
\usepackage{enumitem}

\newcommand{\vect}[1]{\bm{#1}}
\newcommand{\unit}[1]{\,\mathrm{#1}}
\newcommand{\fm}{\unit{fm}}
\newcommand{\GeV}{\unit{GeV}}
\newcommand{\MeV}{\unit{MeV}}

\title{\texttt{genQE\_3N} --- Three-Nucleon GCF Event Generator\\[4pt]
       \large Physics \& Pipeline Documentation}
\author{GCF with 3N Project}
\date{\today}

\begin{document}
\maketitle
\tableofcontents
\newpage

%======================================================================
\section{Overview}
%======================================================================

\texttt{genQE\_3N} extends the two-nucleon GCF generator to
\textbf{three-nucleon short-range correlations (3N SRC)}.  It generates
events of the form
\begin{equation}
  e + A \;\longrightarrow\; e' + N_{1} + N_{2} + N_{3} + (A{-}3)^{*},
\end{equation}
where the electron scatters off one nucleon of a correlated triplet.
The current implementation targets $^{4}$He with $ppn$ triplets.

The key conceptual difference from the 2N generator is that
\emph{contact coefficients times universal functions} are replaced by
\textbf{pre-computed 3D density matrices} $\rho(\theta_{ab},\,k_{\text{cm}},\,k_{\text{rel}})$
obtained from \textit{ab initio} nuclear structure calculations.

%======================================================================
\section{From 2N Contacts to 3N Density Matrices}
%======================================================================

\subsection{2N recap}

In the 2N GCF the pair momentum distribution factorizes as
\begin{equation}
  n^{(ij)}(\vect{k}) = C_{ij}^{A}\,|\varphi_{ij}(k)|^{2},
\end{equation}
separating a nucleus-dependent contact $C_{ij}^{A}$ from a universal function
$\varphi_{ij}$.

\subsection{3N generalization}

For three nucleons no analogous simple factorization exists.  Instead, the
full three-body correlation is encoded in a density matrix that depends on
three independent kinematic variables describing the internal configuration
of the triplet:
\begin{equation}
  \rho_{3N}(\theta_{ab},\,k_{\text{cm}},\,k_{\text{rel}}),
\end{equation}
where
\begin{itemize}[nosep]
  \item $\theta_{ab}$: angle between the two Jacobi-like momenta
        $\vect{p}_{a}$ and $\vect{p}_{b}$ (see below),
  \item $k_{\text{cm}}$: magnitude of the three-body center-of-mass--like
        momentum,
  \item $k_{\text{rel}}$: magnitude of the relative momentum within the pair.
\end{itemize}

These density matrices are tabulated on a regular grid and loaded from
binary data files at initialization.

\subsection{Data files}

Four NN-potential choices are available:
\begin{center}
\begin{tabular}{cl}
\toprule
Index & File \\
\midrule
1 & \texttt{AV8\_1D.dnst}  (Argonne $v_8'$) \\
2 & \texttt{N2LO\_1D.dnst} (Chiral N$^{2}$LO) \\
3 & \texttt{G3\_1D.dnst}   (Chiral N$^{3}$LO / higher) \\
4 & \texttt{AV4\_1D.dnst}  (Argonne AV4$'$) \\
\bottomrule
\end{tabular}
\end{center}

Each file contains $61 \times 101 \times 101 = 622{,}261$ entries:
\begin{center}
\begin{tabular}{lcl}
\toprule
Variable & Bins & Range \\
\midrule
$\theta_{ab}$ & 61  & $[0^{\circ},\,180^{\circ}]$, $\Delta\theta = 3^{\circ}$ \\
$k_{\text{cm}}$  & 101 & $[0,\,10]\fm^{-1}$, $\Delta k = 0.1\fm^{-1}$ \\
$k_{\text{rel}}$ & 101 & $[0,\,10]\fm^{-1}$, $\Delta k = 0.1\fm^{-1}$ \\
\bottomrule
\end{tabular}
\end{center}
Values are trilinearly interpolated at runtime.

%======================================================================
\section{Three-Body Kinematics}
%======================================================================

\subsection{Coordinate system}

The three nucleon momenta are parameterized via two Jacobi-like vectors
$\vect{p}_{a}$, $\vect{p}_{b}$ and a three-body center-of-mass momentum
$\vect{P}_{\text{CM}}$.

\paragraph{Internal frame.}
$\vect{p}_{a}$ is placed along $\hat{z}$ and $\vect{p}_{b}$ lies in the
$xz$-plane:
\begin{equation}
  \vect{p}_{a} = (0,\,0,\,p_{a}), \qquad
  \vect{p}_{b} = (p_{b}\sin\theta_{ab},\;0,\;p_{b}\cos\theta_{ab}).
\end{equation}

\paragraph{Rotation to lab.}
Three Euler angles $(\phi_{a},\,\theta_{a},\,\phi_{baz})$, sampled uniformly,
rotate this frame into the laboratory.

\paragraph{Individual momenta.}
\begin{equation}
  \vect{v}_{1} = \vect{p}_{a},
  \qquad
  \vect{v}_{2} = -\tfrac{1}{2}\,\vect{p}_{a} + \vect{p}_{b},
  \qquad
  \vect{v}_{3} = -\tfrac{1}{2}\,\vect{p}_{a} - \vect{p}_{b}.
\end{equation}
After Euler rotation, each nucleon receives an additional
$\vect{P}_{\text{CM}}/3$ boost.  The $(A{-}3)$ residual recoils with
$\vect{p}_{A{-}3} = -\vect{P}_{\text{CM}}$.

\subsection{Center-of-mass motion}

\begin{equation}
  P_{\text{CM},\alpha} \sim \mathcal{N}(0,\;\sigma_{\text{CM}}),
  \qquad
  \sigma_{\text{CM}} = 0.055\GeV/c.
\end{equation}

%======================================================================
\section{Event Generation Pipeline}
%======================================================================

\subsection{Step 1: Triplet type assignment}

Currently only $ppn$ triplets are generated.  All three nucleons are
initialized as protons, then exactly one --- chosen uniformly at random
among $N_{1}$, $N_{2}$, $N_{3}$ --- is reassigned to a neutron.  Thus
the leading nucleon $N_{1}$ can be either a proton or a neutron (with
probability $2/3$ and $1/3$ respectively).  The weight carries a factor
of~3 to account for the three equivalent assignments.

\subsection{Step 2: Sample three-body CM momentum}

\begin{equation}
  \vect{P}_{\text{CM}} \sim \mathcal{N}^{3}(0,\,\sigma_{\text{CM}}).
\end{equation}

\subsection{Step 3: Sample orientation (Euler angles)}

\begin{equation}
  \phi_{a} \in [-\pi,\,\pi], \qquad
  \cos\theta_{a} \in [-1,\,1], \qquad
  \phi_{baz} \in [-\pi,\,\pi].
\end{equation}
Phase-space weight factor: $\Delta\phi_{a}\,\Delta(\cos\theta_{a})\,\Delta\phi_{baz} = 8\pi^{2}$.

\subsection{Step 4: Sample relative momenta}

\begin{equation}
  p_{a} \in [0,\,5.0]\GeV/c, \qquad
  p_{b} \in [0,\,5.0]\GeV/c, \qquad
  \theta_{ab} \in [10^{-4},\,\pi].
\end{equation}
Phase-space weight factor: $5.0 \times 5.0 \times \pi$.

\subsection{Step 5: Construct individual momenta}

Build $\vect{v}_{1,2,3}$ from the Jacobi decomposition, apply the Euler
rotation, add $\vect{P}_{\text{CM}}/3$ to each, and set
$\vect{p}_{A{-}3} = -\vect{P}_{\text{CM}}$.

\subsection{Step 6: Residual-nucleus energy}

For $^{4}$He the $(A{-}3) = {}^{1}$H system is a single nucleon:
\begin{equation}
  m_{A{-}3} = m_{N} + E^{*},
  \qquad
  E_{A{-}3} = \sqrt{|\vect{p}_{A{-}3}|^{2} + m_{A{-}3}^{2}}.
\end{equation}

Spectator energies:
\begin{equation}
  E_{2} = \sqrt{|\vect{v}_{2}|^{2} + m_{N}^{2}},
  \qquad
  E_{3} = \sqrt{|\vect{v}_{3}|^{2} + m_{N}^{2}}.
\end{equation}

Leading nucleon energy from conservation:
\begin{equation}
  E_{1} = m_{A} - E_{A{-}3} - E_{2} - E_{3}.
\end{equation}

\subsection{Step 7: Electron scattering vertex}

The scattered electron direction is sampled uniformly:
\begin{equation}
  \cos\theta_{k} \in [\cos 40^{\circ},\,\cos 5^{\circ}],
  \qquad
  \phi_{k} \in [-\pi,\,\pi].
\end{equation}

Define the recoil system $\vect{X} = \vect{v}_{2} + \vect{v}_{3} + \vect{p}_{A{-}3}$
and $Z = m_{A} + E_{\text{beam}} - E_{2} - E_{3} - E_{A{-}3}$.  The
scattered electron energy is fixed by 4-momentum conservation:
\begin{equation}
  E_{k} = \frac{Z^{2} - m_{N}^{2} - |\vect{X}|^{2}}
               {2\bigl(Z - \vect{X}\cdot\hat{k}\bigr)}.
\end{equation}

The virtual photon is $\vect{q} = \vect{p}_{\text{beam}} - \vect{k}$,
and the lead nucleon receives the photon:
$\vect{p}_{\text{lead}} = \vect{v}_{1} + \vect{q}$.

\subsection{Step 8: Jacobian}

The energy-conservation $\delta$-function produces a Jacobian:
\begin{equation}
  J_{\delta} = E_{\text{lead}} + \vect{p}_{\text{lead}} \cdot
               (\vect{v}_{1} + \vect{p}_{\text{beam}} + \hat{k}),
  \qquad
  w \;\mathrel{*}=\; \frac{E_{\text{lead}}}{J_{\delta}}.
\end{equation}

\subsection{Step 9: Cross section and density weight}

\begin{equation}
  w \;\mathrel{*}=\; \frac{1}{2}\,(2\pi)^{-7}\,C\,t
    \;\times\; \sigma_{eN}(E_{\text{beam}},\,\vect{k},\,\vect{p}_{\text{lead}})
    \;\times\; \rho_{3N}(\theta_{ab},\,p_{a},\,p_{b}),
\end{equation}
where $C = 0.12$ and $t = 2$ are normalization constants.
$\sigma_{eN}$ uses the same Kelly form factors as the 2N generator.

\subsection{Step 10: Density matrix evaluation}

The sampled variables $(p_{a},\,p_{b},\,\theta_{ab})$ must be
transformed to the density-matrix coordinates
$(k_{\text{cm}},\,k_{\text{rel}},\,\theta'_{ab})$ via a coordinate
relabeling that depends on which pair $(N_{2},N_{3})$ is the
correlated subsystem.  The Jacobian of this coordinate transformation is
\begin{equation}
  J_{\text{coord}}
  = \frac{p_{a}\,p_{b}\,\sin\theta_{ab}}
         {p'_{a}\,p'_{b}\,\sin\theta'_{ab}},
\end{equation}
and the density is
\begin{equation}
  \rho = J_{\text{coord}} \times
         \text{interpolate}\bigl[\rho_{\text{table}}(\theta'_{ab},\,k_{\text{cm}},\,k_{\text{rel}})\bigr].
\end{equation}

%======================================================================
\section{Weight Summary}
%======================================================================

Collecting all factors:
\begin{equation}
  w = 3 \times 8\pi^{2} \times 25\pi
      \times \Delta(\cos\theta_{k})\,\Delta\phi_{k}
      \times \frac{E_{\text{lead}}}{J_{\delta}}
      \times \frac{C\,t}{2\,(2\pi)^{7}}
      \times \sigma_{eN}
      \times \rho_{3N}
      \times \Theta(w > 0).
\end{equation}

%======================================================================
\section{Alternative Momentum Parameterization}
%======================================================================

For analysis purposes (Dalitz/ternary diagrams, spectral functions), an
alternative parameterization in terms of momentum fractions is available:
\begin{equation}
  f_{i} = \frac{|\vect{p}_{i}|}{|\vect{p}_{1}|+|\vect{p}_{2}|+|\vect{p}_{3}|},
  \qquad
  f_{1} + f_{2} + f_{3} = 1.
\end{equation}
The Jacobian for the transformation
$(\vect{k}_{1},\vect{k}_{2},\vect{k}_{3}) \to (p_{\text{tot}},f_{1},f_{2})$
is
\begin{equation}
  J = \frac{8}{\sqrt{3}}\,\frac{1}{p_{1}\,p_{2}\,p_{3}\,p_{\text{tot}}^{2}}.
\end{equation}

%======================================================================
\section{Differences from the 2N Generator}
%======================================================================

\begin{center}
\begin{tabular}{lll}
\toprule
Aspect & \texttt{genQE} (2N) & \texttt{genQE\_3N} (3N) \\
\midrule
Correlated cluster & Pair ($N N$) & Triplet ($N N N$) \\
Structure input    & Contacts $C_{ij}\,\varphi^{2}(k)$ & 3D density matrix $\rho(\theta,k_{\text{cm}},k_{\text{rel}})$ \\
Residual system    & $(A{-}2)^{*}$ & $(A{-}3)^{*}$ \\
Phase space dim.   & 6D (CM + rel) & 9D (CM + Euler + Jacobi) \\
Electron vertex    & Quadratic solve for $k_{\perp}$ & Direct solve for $E_{k}$ \\
$\sigma_{\text{CM}}$ & 0.10--0.15 GeV/$c$ & 0.055 GeV/$c$ \\
Target nuclei      & $^{2}$H through $^{208}$Pb & $^{4}$He \\
\bottomrule
\end{tabular}
\end{center}

%======================================================================
\section{Summary of Algorithm}
%======================================================================

For each Monte Carlo trial:
\begin{enumerate}
  \item Assign triplet types ($ppn$); weight $\times\,3$.
  \item Sample $\vect{P}_{\text{CM}}$ (Gaussian, $\sigma = 0.055\GeV/c$).
  \item Sample Euler angles $(\phi_{a},\theta_{a},\phi_{baz})$ uniformly.
  \item Sample Jacobi magnitudes $(p_{a},p_{b})$ and angle $\theta_{ab}$
        uniformly.
  \item Construct $\vect{v}_{1,2,3}$ via Jacobi decomposition + Euler
        rotation + CM boost.
  \item Compute residual and spectator energies; derive $E_{1}$ from
        conservation.
  \item Sample electron direction $(\theta_{k},\phi_{k})$; solve for $E_{k}$.
  \item Compute Jacobian $J_{\delta}$ from energy-conservation constraint.
  \item Evaluate $\sigma_{eN}$ (CC1, Kelly form factors).
  \item Look up $\rho_{3N}$ from density matrix with coordinate transform.
  \item Multiply all weight factors; store event if $w > 0$.
\end{enumerate}

\end{document}
